\chapter{Demande de cotation et marchés de dealers en théorie}\label{ch:mx:request-for-quote-and-dealer-markets-in-theory}

Un client qui demande un prix à cinq dealers pour une obligation d'entreprise obtient un meilleur prix que celui qui en interroge deux, jusqu'à un certain point~: chaque dealer sollicité apprend
qu'un vendeur existe, et ceux qui perdent la demande peuvent trader sur cette information. Ce chapitre présente la théorie des marchés qui n'ont pas de carnet d'ordres central~: un modèle de recherche et de négociation des
fourchettes des dealers, la demande de cotation (RFQ) vue comme une enchère à valeurs privées, à valeurs communes et avec \omterm{def:mx:dark-pools-and-blocks:leakage}{fuite d'information}, le
nombre de dealers qu'un client doit interroger, et les réseaux et relations par lesquels les dealers échangent.

\section{Recherche et négociation}

\begin{definition}[Modèle de recherche de gré à gré, friction de recherche, pouvoir de négociation]\label{def:mx:request-for-quote-and-dealer-markets-in-theory:search}
Un \emph{modèle de recherche de gré à gré}\index{modèle de recherche de gré à gré} décrit un marché dans lequel les traders trouvent leurs contreparties une à une,
à des instants aléatoires, et s'accordent sur les prix de gré à gré. Une \emph{friction de recherche}\index{friction de recherche} est le délai et le coût nécessaires pour trouver une contrepartie~; le
\emph{pouvoir de négociation}\index{pouvoir de négociation} d'un trader est la part des gains d'une transaction qu'il capte lorsqu'une contrepartie est trouvée.
\end{definition}

Duffie, Gârleanu et Pedersen (2005) ont calculé les prix auxquels les investisseurs traitent avec les dealers dans un modèle dynamique de recherche et de négociation, et ont montré
que les fourchettes sont plus étroites lorsque les investisseurs trouvent plus facilement d'autres investisseurs ou atteignent plusieurs dealers. Ce chapitre résout une version de leur
marché de dealers (\texttt{firm\_otcsearch.dgp}). Un actif verse une unité par an~; les investisseurs sont de type élevé ou de type faible, le type faible supportant un coût de détention
$\delta$~; les types changent aux taux $\lambda_u$ (vers élevé) et $\lambda_d$ (vers faible). Les propriétaires de type faible vendent aux dealers et les non-propriétaires de type élevé achètent auprès d'eux,
chaque investisseur rencontrant un dealer au taux $\rho$ (un processus de Poisson, livre~4, chapitre~6)~; les dealers couvrent chaque transaction sur un marché inter-dealers et conservent une
part $z$ du surplus de chaque transaction. Avec une offre $s$ inférieure à la part des types élevés $\lambda_u/(\lambda_u+\lambda_d)$, les acheteurs sont le côté long et le
prix inter-dealers est leur valeur de réserve.

En régime stationnaire, les propriétaires de type faible affluent par changement de type et sortent vers les dealers et vers le type élevé~: $\lambda_d\mu_{ho}=(\lambda_u+\rho)\mu_{lo}$,
avec $\mu_{ho}+\mu_{lo}=s$. Les valeurs de réserve $\Delta V_\sigma$ (la valeur de la détention par rapport à la non-détention, pour le type $\sigma$) sont solutions de deux équations
linéaires (les équations de Bellman d'une chaîne de Markov à temps continu, livre~4, chapitre~8)~:
\[
r\Delta V_h=1+\lambda_d(\Delta V_l-\Delta V_h),\qquad r\Delta V_l=1-\delta+\lambda_u(\Delta V_h-\Delta V_l)+\rho(1-z)(\Delta V_h-\Delta V_l),
\]
et le cours vendeur est $\Delta V_h$, le cours acheteur $z\Delta V_l+(1-z)\Delta V_h$, la fourchette $z(\Delta V_h-\Delta V_l)$.

\omcode{code/firm/otcsearch/firm_otcsearch.py}{43}{58}{Le régime stationnaire du marché de dealers et ses prix~: les masses de propriétaires et de non-propriétaires de chaque type, puis les deux valeurs de réserve issues de leurs équations de Bellman linéaires, et le cours acheteur et le cours vendeur que produit la négociation.}\label{lst:mx:request-for-quote-and-dealer-markets-in-theory:dgp}

Avec $r=5~\%$, $\lambda_u=2$ et $\lambda_d=0.2$ par an, $s=0.8$, $\delta=2$ et $z=0.8$, la fourchette passe de 210 points de base du prix lorsque
les investisseurs rencontrent un dealer dix fois par an, à 36.7 points de base à cent rencontres et à 4.0 à mille (\cref{fig:mx:request-for-quote-and-dealer-markets-in-theory:dgp})~:
une recherche plus rapide laisse aux types faibles moins de temps pour supporter leur coût de détention, et c'est ce surplus que les dealers négocient.

\begin{omfigure}
\begin{tikzpicture}
\begin{axis}[width=9cm,height=5.4cm,xmode=log,ymode=log,xlabel={\small rencontres avec les dealers par an ($\rho$)},ylabel={\small fourchette (pb du prix)},
  tick label style={font=\footnotesize},grid=major,grid style={gray!25},table/col sep=comma]
\addplot[omDef,very thick,mark=*] table[x=rho,y=spread_bp]{figdata/microstructure/23-request-for-quote-and-dealer-markets-in-theory/dgp.csv};
\end{axis}
\end{tikzpicture}

\omcaption{La fourchette acheteur--vendeur du marché de dealers en fonction du taux de contact des investisseurs avec les dealers, dans le modèle de recherche et de négociation avec les
paramètres du chapitre. Données~: \texttt{mx\_otc.dgp\_study}.}\label{fig:mx:request-for-quote-and-dealer-markets-in-theory:dgp}
\end{omfigure}

\section{Réseaux de dealers}

\begin{definition}[Réseau cœur--périphérie]\label{def:mx:request-for-quote-and-dealer-markets-in-theory:cp}
Un \emph{réseau cœur--périphérie}\index{réseau cœur--périphérie} de dealers comprend un petit cœur de dealers qui traitent entre eux et avec de nombreux autres,
et une large périphérie dont les dealers traitent surtout avec le cœur.
\end{definition}

Li et Schürhoff (2019) ont constaté que les dealers d'obligations municipales forment un tel réseau, avec 10 à 30 plaques tournantes et plus de 2\,000 broker-dealers périphériques~; les obligations circulent
de la périphérie vers le cœur et, en partie, en sens inverse~; les dealers centraux facturent aux investisseurs jusqu'au double des marges aller-retour des dealers périphériques, mais assurent l'immédiateté en
appariant plus directement acheteurs et vendeurs. Le réseau modèle du chapitre (\texttt{core\_periphery}) compte 10 dealers de cœur et 200 dealers périphériques
reliés chacun à deux dealers de cœur~; une obligation emprunte le chemin le plus court du dealer du vendeur à celui de l'acheteur, et chaque dealer sur le chemin
prélève une marge, 4 points de base au cœur et 2 en périphérie. Un client d'un dealer du cœur est servi en 1.77 étapes en moyenne et paie 9.2 points de
base~; un client d'un dealer périphérique, 2.58 étapes et 10.4. La longueur de la chaîne fait partie du prix~; le constat empirique selon lequel le cœur facture davantage
provient de ses marges, que ce modèle fixe à la main.

\section{La demande de cotation comme enchère}

Une demande de cotation (livre~2, chapitre~22) est une enchère au premier prix sous pli scellé entre les dealers interrogés (livre~4, chapitre~29). Chaque
dealer attribue à l'obligation une valeur qui comporte une part privée (son inventaire, ses axes~: écart-type $\sigma_p$), et il estime la valeur commune avec une
erreur ($\sigma_c$). Le client vend au meilleur cours acheteur, et chaque dealer perdant, qui sait désormais qu'un vendeur est sur le marché, coûte en moyenne au client une \omterm{def:mx:dark-pools-and-blocks:leakage}{fuite d'information} $\ell$
en moyenne. Un dealer qui remporte l'enchère alors que son estimation de la valeur commune est la plus élevée l'a surestimée~: il minore son offre de $\sigma_c\E[\max\text{ de
}n]$ (la malédiction du vainqueur, livre~2, chapitre~30, et les \omterm{def:mx:portfolio-trades-and-risk-bids:agency}{risk bids} du chapitre~22). Le coût espéré du client, en dessous de la valeur commune, est
\[
C(n)=m-\bigl(\sqrt{\sigma_c^2+\sigma_p^2}-\sigma_c\bigr)\E[\max\text{ de }n]+\ell(n-1),
\]
avec $m$ la marge des dealers~: la concurrence aide grâce aux valeurs privées, la partie commune se compense en grande partie avec la minoration, et chaque dealer supplémentaire
ajoute une \omterm{def:mx:dark-pools-and-blocks:leakage}{fuite d'information} (\texttt{rfq\_cost}, sur \texttt{firm.rfq} du livre~2~; une simulation de l'enchère concorde à 0.001 point de base près).

\omcode{code/firm/otcsearch/firm_otcsearch.py}{61}{70}{Le coût espéré d'une demande de cotation adressée à n dealers, construit à partir de la formule à valeurs privées du livre~2, et le nombre de dealers qui le minimise.}\label{lst:mx:request-for-quote-and-dealer-markets-in-theory:rfq}

\section{Combien de dealers interroger}

Avec une marge de 10 points de base, $\sigma_p=5$ et $\sigma_c=2$, la \cref{fig:mx:request-for-quote-and-dealer-markets-in-theory:rfq} trace $C(n)$ pour
quatre coûts de fuite, et le \cref{tab:mx:request-for-quote-and-dealer-markets-in-theory:bestn} donne l'optimum.

\begin{omfigure}
\begin{tikzpicture}
\begin{axis}[width=10cm,height=5.8cm,xlabel={\small dealers interrogés ($n$)},ylabel={\small coût espéré (pb)},
  tick label style={font=\footnotesize},grid=major,grid style={gray!25},xmin=0.5,xmax=12.5,ymin=5,ymax=16,
  legend columns=4,legend style={font=\footnotesize,at={(0.5,-0.3)},anchor=north,draw=none,fill=none,
  /tikz/every even column/.append style={column sep=6pt}},table/col sep=comma]
\addplot[omDef,thick,mark=*,mark size=1.3pt] table[x=n,y=l025]{figdata/microstructure/23-request-for-quote-and-dealer-markets-in-theory/rfq.csv};
\addplot[omThm,thick,mark=square*,mark size=1.3pt] table[x=n,y=l05]{figdata/microstructure/23-request-for-quote-and-dealer-markets-in-theory/rfq.csv};
\addplot[omProp,thick,mark=triangle*,mark size=1.6pt] table[x=n,y=l1]{figdata/microstructure/23-request-for-quote-and-dealer-markets-in-theory/rfq.csv};
\addplot[gray,thick,mark=diamond*,mark size=1.6pt] table[x=n,y=l2]{figdata/microstructure/23-request-for-quote-and-dealer-markets-in-theory/rfq.csv};
\legend{{$\ell=0.25$},{$\ell=0.5$},{$\ell=1$},{$\ell=2$}}
\end{axis}
\end{tikzpicture}

\omcaption{Le coût espéré de la vente par demande de cotation en fonction du nombre de dealers interrogés, pour quatre coûts de fuite par dealer perdant
(points de base)~: marge 10, dispersion des valeurs privées 5, erreur sur la valeur commune 2. Données~: \texttt{mx\_otc.rfq\_study}.}\label{fig:mx:request-for-quote-and-dealer-markets-in-theory:rfq}
\end{omfigure}

\begin{table}[ht]
\centering
\small
\begin{tabular}{@{}lrrrrrr@{}}
\toprule
fuite par dealer perdant (pb) & 0.1 & 0.25 & 0.5 & 1 & 2 & 3 \\
\midrule
dealers à interroger & 16 & 7 & 4 & 2 & 1 & 1 \\
coût espéré (pb) & 5.52 & 6.92 & 8.02 & 9.09 & 10.00 & 10.00 \\
\bottomrule
\end{tabular}
\caption{Le nombre de dealers qui minimise le coût et le coût espéré en fonction du coût de fuite par dealer perdant. Données~: \texttt{mx\_otc.rfq\_study}.}\label{tab:mx:request-for-quote-and-dealer-markets-in-theory:bestn}
\end{table}

L'optimum décroît vite~: de 16 dealers lorsqu'un perdant coûte un dixième de point de base à 4 pour un demi-point de base, et à un seul dealer à partir de 2
points de base. La part de valeur commune compte aussi~: pour un demi-point de base de fuite, un client devrait interroger 6 dealers si les estimations des dealers n'avaient aucune
erreur commune, 4 avec $\sigma_c=2$ et 3 avec $\sigma_c=5$, car la part commune ajoute du bruit que le gagnant réduit sans ajouter de concurrence.
Riggs, Onur, Reiffen et Zhu (2020) ont constaté que les clients des deux plus grandes plateformes d'exécution de swaps dealer-client sur CDS d'indices contactent
en général peu de dealers, et ont modélisé la raison~: le bénéfice d'une exposition plus large est atténué par la malédiction du vainqueur et par les relations client--dealer.

\section{Relations et chasse à l'information}

\begin{definition}[Chasse à l'information]\label{def:mx:request-for-quote-and-dealer-markets-in-theory:chasing}
La \emph{chasse à l'information}\index{chasse à l'information} est le fait pour un dealer de coter de meilleurs prix à des clients dont les transactions portent de l'information, car
ce qu'il apprend de leurs flux vaut davantage pour lui, dans ses transactions avec d'autres, que la sélection adverse qu'il subit sur leurs transactions.
\end{definition}

Dans une cotation concurrentielle, la fourchette d'un dealer envers un client est son coût plus la sélection adverse qu'il anticipe, diminué de la valeur de ce qu'il apprend. Avec un
coût de 2 points de base et une sélection adverse de 1.5 sur la transaction d'un client informé, ce client est coté 3.5 lorsque son information
ne vaut rien pour le dealer, la même 2.0 qu'un client non informé lorsqu'elle vaut 1.5, et 0.5 lorsqu'elle vaut 3 (\texttt{chase})~: le client informé
obtient le meilleur prix. Les relations fonctionnent de la même façon~: Hendershott, Li, Livdan et Schürhoff (2020) ont constaté, sur les transactions d'obligations d'entreprise d'assureurs, que
les prix d'exécution s'améliorent d'abord avec le nombre de dealers dans le réseau d'un client, puis se dégradent au-delà d'une vingtaine, car les affaires répétées avec quelques dealers se compensent
avec la concurrence entre un grand nombre~; Di Maggio, Kermani et Song (2017) ont constaté que les dealers appliquent des fourchettes plus faibles aux dealers avec lesquels ils entretiennent les liens les plus forts, davantage en période de turbulences, et que
les dealers d'importance systémique facturent davantage aux dealers périphériques et aux clients.

\section{Tutoriel~: dealers, recherches et enchères}\label{tut:mx:request-for-quote-and-dealer-markets-in-theory:tut}

\begin{tutorial}
\textbf{Objectif.} Résoudre un marché de dealers à recherche et négociation, trouver le nombre de dealers minimisant le coût dans une demande de cotation, et tracer
les chaînes d'intermédiation dans un \omterm{def:mx:request-for-quote-and-dealer-markets-in-theory:cp}{réseau cœur--périphérie}.
\textbf{État final~:} \cref{fig:mx:request-for-quote-and-dealer-markets-in-theory:dgp,fig:mx:request-for-quote-and-dealer-markets-in-theory:rfq},
le \cref{tab:mx:request-for-quote-and-dealer-markets-in-theory:bestn} et les nombres des sections 2 et 5.
\begin{tutsteps}
\item \textbf{Recherche.} \texttt{firm\_otcsearch.dgp(lam\_u, lam\_d, rho, s, r, delta, z)}~; \texttt{mx\_otc.dgp\_study()}.
\item \textbf{Enchère.} \texttt{rfq\_cost}, \texttt{best\_n}, \texttt{rfq\_simulate}~; \texttt{rfq\_study()}, \texttt{common\_study()}.
\item \textbf{Réseaux et chasse.} \texttt{core\_periphery}, \texttt{intermediation}, \texttt{chase}~; \texttt{network\_study()},
\texttt{chase\_study()}~; tracer avec \texttt{fig\_otc.py}.
\end{tutsteps}
\textbf{Pour aller plus loin.} Laisser aussi les investisseurs se rencontrer directement (l'autre canal de Duffie, Gârleanu et Pedersen)~; faire dépendre la fuite des
inventaires des perdants~; laisser les dealers périphériques choisir leurs liens avec le cœur.
\end{tutorial}

\section{Construction~: recherche OTC et enchères}\label{bld:mx:request-for-quote-and-dealer-markets-in-theory:otcsearch}

\begin{build}
\bfield{Rôle} La théorie derrière l'adaptateur obligataire du chapitre~21 et les choix de dealers du desk~: prix issus de la recherche, nombre de dealers adéquat et
chaînes d'intermédiation.
\bfield{Interface} \texttt{dgp(lam\_u, lam\_d, rho, s, r, delta, z)}, \texttt{rfq\_cost(n, markup, sigma\_p, sigma\_c, leak)}, \texttt{best\_n},
\texttt{rfq\_simulate}, \texttt{chase(cost, adverse, info\_value)}, \texttt{core\_periphery(n\_core, n\_periph, links, seed)},
\texttt{intermediation(adj, n\_core, markup\_core, markup\_periph, trials, seed)}.
\bfield{Règles} L'actif verse un euro par an~; les coûts sont exprimés en points de base sous la valeur commune~; la version du modèle de recherche avec des acheteurs du côté long.
\bfield{Tests d'acceptation} \texttt{code/firm/otcsearch/tests/}~: le bilan des flux et les masses de l'état stationnaire~; la fourchette diminue avec le taux de contact et
s'annule sans \omterm{def:mx:request-for-quote-and-dealer-markets-in-theory:search}{pouvoir de négociation} des dealers~; la formule de la RFQ comparée à celle du livre~2 et à la simulation~; la monotonie de l'optimum en fonction de la fuite~; la chasse
réduit la fourchette informée~; les transactions périphériques ont des chaînes plus longues.
\bfield{Pour aller plus loin} Recherche d'investisseur à investisseur~; entrée endogène des dealers~; fuite dépendant des inventaires.
\end{build}

\begin{omsources}
\item D.~Duffie, N.~Gârleanu et L.~H.~Pedersen, ``Over-the-counter markets'', \emph{Econometrica} 73(6), 2005.
\item P.~Di~Maggio, A.~Kermani et Z.~Song, ``The value of trading relations in turbulent times'', \emph{Journal of Financial Economics} 124(2),
2017.
\item D.~Li et N.~Schürhoff, ``Dealer networks'', \emph{Journal of Finance} 74(1), 2019.
\item L.~Riggs, E.~Onur, D.~Reiffen et H.~Zhu, ``Swap trading after Dodd--Frank: evidence from index CDS'', \emph{Journal of Financial Economics}
137(3), 2020.
\item T.~Hendershott, D.~Li, D.~Livdan et N.~Schürhoff, ``Relationship trading in over-the-counter markets'', \emph{Journal of Finance} 75(2), 2020.
\end{omsources}

\section{Exercices}

\begin{exercise}[$\star$]\label{exo:mx:request-for-quote-and-dealer-markets-in-theory:1}
Avec $\lambda_u=2$, $\lambda_d=0.2$, $s=0.8$ et $\rho=100$, calculer la masse des propriétaires de type faible dans l'état stationnaire.
\end{exercise}

\begin{exercise}[$\star$]\label{exo:mx:request-for-quote-and-dealer-markets-in-theory:2}
Calculer $C(n)$ pour $n=1$ et $n=2$ avec les paramètres du chapitre et une fuite de 1 point de base.
\end{exercise}

\begin{exercise}[$\star$]\label{exo:mx:request-for-quote-and-dealer-markets-in-theory:3}
Pour quelle valeur de l'information d'une transaction informée un dealer cote-t-il la même fourchette à des clients informés et non informés~?
\end{exercise}

\begin{exercise}[$\star\star$]\label{exo:mx:request-for-quote-and-dealer-markets-in-theory:4}
Pourquoi la fourchette s'annule-t-elle lorsque les dealers n'ont aucun \omterm{def:mx:request-for-quote-and-dealer-markets-in-theory:search}{pouvoir de négociation}, et pourquoi diminue-t-elle avec le taux de contact~?
\end{exercise}

\begin{exercise}[$\star\star$]\label{exo:mx:request-for-quote-and-dealer-markets-in-theory:5}
Pourquoi une erreur de valeur commune réduit-elle le nombre de dealers qu'un client doit interroger~?
\end{exercise}

\begin{exercise}[$\star\star$]\label{exo:mx:request-for-quote-and-dealer-markets-in-theory:6}
Dans le réseau jouet, pourquoi les clients des dealers périphériques passent-ils par des chaînes plus longues, et qu'est-ce qui leur ferait payer moins~?
\end{exercise}

\begin{exercise}[$\star\star\star$]\label{exo:mx:request-for-quote-and-dealer-markets-in-theory:7}
\emph{Programmation.} Sans erreur de valeur commune et avec une fuite de 1 point de base, combien de dealers le client doit-il interroger, et combien cela coûte-t-il~?
\end{exercise}

\begin{exercise}[$\star\star\star$]\label{exo:mx:request-for-quote-and-dealer-markets-in-theory:8}
\emph{Trouver la faille.} «~Davantage de dealers signifie toujours davantage de concurrence, donc nos RFQ partent vers chaque dealer de la plateforme.~»
\end{exercise}

\section{Problème~: combien de dealers interroger}

\begin{problem}[{Problème du week-end --- combien de dealers interroger}]\label{pb:mx:request-for-quote-and-dealer-markets-in-theory:1}
Le desk obligataire d'un fonds envoie chaque RFQ à huit dealers. Le responsable du trading demande si c'est le bon choix.

\medskip\noindent\textbf{Partie I --- Recherche.}
\begin{enumerate}
\item Définir une \omterm{def:mx:request-for-quote-and-dealer-markets-in-theory:search}{friction de recherche} et le \omterm{def:mx:request-for-quote-and-dealer-markets-in-theory:search}{pouvoir de négociation}.
\item Énoncer les types, les taux et les transactions du modèle de recherche.
\item Dériver la masse de propriétaires de type faible dans l'état stationnaire.
\item Écrire les équations de Bellman pour les valeurs de réserve et pour le cours acheteur et le cours vendeur.
\item Comment la fourchette dépend-elle du taux de contact~?
\end{enumerate}

\medskip\noindent\textbf{Partie II --- Les enchères.}
\begin{enumerate}[resume]
\item Pourquoi une RFQ est-elle une enchère au premier prix sous pli scellé~?
\item Dériver la minoration de l'enchérisseur due à la malédiction du vainqueur pour une erreur à valeur commune.
\item Écrire le coût espéré du client $C(n)$ et expliquer chaque terme.
\item Dans quelle mesure une simulation de l'enchère concorde-t-elle avec le modèle~?
\end{enumerate}

\medskip\noindent\textbf{Partie III --- Le nombre de dealers.}
\begin{enumerate}[resume]
\item \emph{Énoncer le résultat nommé}~: le nombre de dealers qui minimise le coût et sa dépendance au coût de fuite par dealer perdant.
\item Comment l'erreur à valeur commune le modifie-t-elle~?
\item Qu'ont établi Riggs, Onur, Reiffen et Zhu sur le nombre de dealers que les clients contactent, et pourquoi~?
\item Huit dealers est-il le bon choix pour ce desk~? Que faudrait-il savoir~?
\end{enumerate}

\medskip\noindent\textbf{Partie IV --- Réseaux et relations.}
\begin{enumerate}[resume]
\item Définir un \omterm{def:mx:request-for-quote-and-dealer-markets-in-theory:cp}{réseau cœur--périphérie} et donner les résultats de Li et Schürhoff.
\item Que montre le réseau jouet concernant les chaînes~?
\item Définir la \omterm{def:mx:request-for-quote-and-dealer-markets-in-theory:chasing}{chasse à l'information} et donner les fourchettes qu'elle implique.
\item Qu'ont établi Hendershott, Li, Livdan et Schürhoff sur la taille du réseau et les prix~?
\item Qu'ont établi Di Maggio, Kermani et Song sur les liens entre dealers en période de turbulence~?
\item Comment mesurer le coût de fuite par dealer perdant à partir des données propres au desk~?
\item En une phrase~: que arbitre un client lorsqu'il interroge un dealer de plus~?
\end{enumerate}
\end{problem}

\section{Questions d'entretien}

\begin{interviewq}[$\star$ \iqroles{trader}]\label{iq:mx:request-for-quote-and-dealer-markets-in-theory:1}
Vous devez vendre une obligation d'entreprise peu liquide. Combien de dealers interrogez-vous, et pourquoi~?
\end{interviewq}

\begin{interviewq}[$\star\star$ \iqroles{chercheur}]\label{iq:mx:request-for-quote-and-dealer-markets-in-theory:2}
Expliquer la malédiction du vainqueur dans une RFQ et comment un dealer doit enchérir.
\end{interviewq}

\begin{interviewq}[$\star\star$ \iqroles{chercheur}]\label{iq:mx:request-for-quote-and-dealer-markets-in-theory:3}
Dans un modèle de recherche, pourquoi les fourchettes sont-elles plus faibles lorsque les investisseurs trouvent plus facilement des dealers~?
\end{interviewq}

\begin{interviewq}[$\star\star$ \iqroles{banque}]\label{iq:mx:request-for-quote-and-dealer-markets-in-theory:4}
Pourquoi un dealer pourrait-il coter un prix plus serré à un client qu'il sait informé~?
\end{interviewq}

\begin{interviewq}[$\star\star$ \iqroles{développeur}]\label{iq:mx:request-for-quote-and-dealer-markets-in-theory:5}
Concevoir les données qu'une plateforme RFQ doit conserver pour permettre aux clients de mesurer la fuite.
\end{interviewq}

\begin{interviewq}[$\star\star\star$ \iqroles{risque}]\label{iq:mx:request-for-quote-and-dealer-markets-in-theory:6}
Les dealers centraux d'un marché obligataire subissent un choc. Que deviennent les chaînes d'intermédiation et les coûts supportés par les clients~?
\end{interviewq}
